\documentclass[12pt]{article}
\usepackage[utf8]{inputenc}
\usepackage[T1]{fontenc}
\usepackage{hyperref}
\title{Penang Sign Language Language Resources for K-12 Curriculum Development}
\author{Compiled by Claude Sonnet 5.5 (claude-sonnet-5-5) for Andr\'as Kornai}
\date{}
\begin{document}
\maketitle

\section{Introduction}
Penang Sign Language (ISO 639-3 psg) is a deaf community sign language that arose among deaf children at the Federation School for the Deaf in Penang (founded 1954) and is now used mainly by older signers, most of whom also know Malaysian Sign Language. Search results found no dedicated educational materials, only general descriptions.

\section{Primary Educational Resources}
No textbooks, primers or curricula in Penang Sign Language were found.

\section{Linguistic and Reference Materials}
\begin{itemize}
\item Wikipedia, Penang Sign Language. \url{https://en.wikipedia.org/wiki/Penang_Sign_Language}
\item Wikipedia ISO 639:psg page. \url{https://wikipedia.com/wiki/ISO_639:psg}
\item Rosetta Project (Long Now) entry for psg. \url{https://rosettapanglossia.longnow.org/index.php?title=Psg}
\item A journal article PDF surfaced in the search (content not verified). \url{https://thaiscience.info/journals/Article/CMUS/10325206.pdf}
\end{itemize}

\section{Religious and Community Resources}
\begin{itemize}
\item Penang Deaf Association, mentioned as a partner in sign-language lessons for Starbucks staff. \url{https://thesmartlocal.my/starbucks-penang/amp/}
\item MIND (Malaysian Institute for the Deaf) article on learning Malaysian Sign Language (BIM), relevant to the national sign language used alongside psg. \url{https://mind.org.my/article/learning-malaysian-sign-languagebim/}
\item Joshua Project, Deaf in Malaysia. \url{https://www.joshuaproject.net/index.php/people_groups/19007/MY}
\end{itemize}

\section{Digital and Audio Resources}
No video corpus, dictionary or app for Penang Sign Language was found.

\section{Curriculum Development Context}
The language is endangered and held by older signers; younger deaf people mostly use Malaysian Sign Language. Any curriculum would need documentation work with the Penang deaf community first. Malaysian Sign Language materials may serve as a bridge but are a different language.
\end{document}